% Nouveau EGA V (cote 25): archive cote 25, site page 19 (PDF page 20).
% typescript. Draft compared: none.
% AI-transcribed by Claude Opus 5.5 (claude-opus-5-5, adaptive thinking, effort high), Message Batches API, batch ega25b (msgbatch_01A7n5YuXubUPNK1zsepi8RG), 2026-11-13; not yet checked by eye.
\scanpage{25/19}
\mspage{16}

sections hyperplanes de X lui-même (sans extension préalable du corps
\var{\margin{\unsure{trait vertical et flèche}}}{\placedmargin{\unsure{trait vertical et flèche}}}{marginal note placed}{ega25_25_019-v0001}de base). Cela résulte \unsure{en e ffet|en effet} de la constructibilité de E et du fait
que $\check{P}$ est un schéma de Jacobson.

Donnons à titre d'exemples quelques cas particuliers où les
considérations précédentes s'appliquent:

\uline{Proposition} 4.2. Soit G un Module de présentation finie sur X .
Soit P une des propriétés suivantes pour un Module\var{\inserted{M}}{M}{insertion}{ega25_25_019-v0002} sur un préschéma
algébrique Z sur un corps K ;
\var{\xout{\unsure{xxxthérixxxxxxx(i)xxxxxx(v)xxxxxxxxx2,8xxxAlxxxx}}}{}{typewriter strike-out}{ega25_25_019-v0003}
\var{\deleted{a) L'ensemble \var{\xout{des}}{}{typewriter strike-out}{ega25_25_019-v0005} E des \unsure{$\check{P}$} tels que G satisfasse P est
constructible}}{}{deletion}{ega25_25_019-v0004}
(i) coppof (M) $\unsure{\leq}$ n.

(ii) M satisfait (S$_{k}$).

(i\unsure{\var{\xout{x}}{}{typewriter strike-out}{ega25_25_019-v0006}}i) M est Cohen-Macaulay
\var{\inserted[typed]{cycles}}{cycles}{insertion}{ega25_25_019-v0007}
(iv) M est sans \var{\xout{idéaux}}{}{typewriter strike-out}{ega25_25_019-v0008} premiers associés immergés.

(v) M est séparable sur \var{\xout{\illegible{1}}}{}{typewriter strike-out}{ega25_25_019-v0009} K .

\var{\deleted{(vi) M est géométriquement intègre sur K/}}{}{deletion}{ega25_25_019-v0010}

\var{\xout{(vii)xxMxestxgéomxtxxxintègrexsurxKx}}{}{typewriter strike-out}{ega25_25_019-v0011} Avec ces notations; si E désig%
ne l'ensemble des \hand{$\xi\in$} P tels que G\hand{$_{\xi}$} satisfasse à P, on a :
a) E est une partie constructible de \var{\over{P}{$\check{P}$}}{$\check{P}$}{writing over}{ega25_25_019-v0012} .

b) Supposons S=Spec(k), k un corps, et que F satisfasse à P. Suppo-%
sons de plus \var{\inserted{(dans le cas (v))}}{(dans le cas (v))}{insertion}{ega25_25_019-v0013} que k soit de car 0, ou que f:X$\rightarrow$P soit une \var{\deleted{immersion}}{}{deletion}{ega25_25_019-v0014}\var{\inserted{\unsure{non ramifié}}}{\unsure{non ramifié}}{insertion}{ega25_25_019-v0015}. \var{\xout{\unsure{Akx}}}{}{typewriter strike-out}{ega25_25_019-v0016}
Alors E contient un ouvert dense de \var{\over{P}{$\check{P}$}}{$\check{P}$}{writing over}{ega25_25_019-v0017} .

Démonstration. a) résulte du fait que P est une propriété
constructible, ce qui a été vu au par.9. Quant à b), cela résulte
des résultats correspondants des \var{\xout{deux}}{}{typewriter strike-out}{ega25_25_019-v0018}\var{\inserted{deux}}{deux}{insertion}{ega25_25_019-v0019} n$^{o}$s précédents. Remords\var{\inserted{à b)}}{à b)}{insertion}{ega25_25_019-v0020}: suppo-%
ser plus généralement que si Z est l'ensemble des points de X enlesquel%
F ne satisfait pas à P, on ait f(Z) fini i.e. dim $\overline{f(Z)}\;\unsure{\leq}0$.

\mspage{19}
